% =====================================================================
%  CHƯƠNG 2. SỰ TIẾN HÓA MÔ HÌNH: TỪ CHATBOT ĐẾN AGENTIC MESH
% =====================================================================
\chapter{Sự tiến hóa mô hình: từ Chatbot đến Agentic Mesh}

\section{Đối sánh Non-agentic Chatbot, AI Copilot và Autonomous Agent}
Để làm rõ khác biệt về bản chất kỹ thuật, Bảng~\ref{tab:ch2-compare} đối sánh ba thế hệ công nghệ theo các tiêu chí vận
hành tổng hợp từ Google Cloud \cite{google_agents} và McKinsey \cite{mckinsey_agents}. Google Cloud phân biệt rõ
\emph{bot} (làm theo quy tắc định sẵn, ít tự chủ), \emph{trợ lý AI} (hỗ trợ người dùng, người dùng vẫn ra quyết định) và
\emph{tác tử} (tự chủ cao nhất, tự đưa ra và thực hiện quyết định để đạt mục tiêu).

\begin{table}[!htbp]\centering
\caption{Đối sánh Chatbot truyền thống, trợ lý AI Copilot và hệ tác tử tự trị.}\label{tab:ch2-compare}
\small
\begin{tabularx}{\linewidth}{L{2.9cm}XXX}
\toprule
\textbf{Tiêu chí} & \textbf{Non-agentic Chatbot} & \textbf{AI Copilot (trợ lý)} & \textbf{Autonomous AI Agent} \\
\midrule
Mức độ tự trị & Thụ động, phụ thuộc hoàn toàn vào câu lệnh & Hỗ trợ từng bước, người dùng quyết định & Tự trị cao,
tự điều hướng tới mục tiêu \\
Lập kế hoạch & Không có & Con người tự lập kế hoạch & Tự phân rã nhiệm vụ nhiều bước \\
Sử dụng công cụ & Không có hoặc rất hạn chế & Gợi ý, con người xác nhận từng lệnh & Tự lựa chọn và thực thi công cụ \\
Quản lý bộ nhớ & Ngữ cảnh phiên ngắn & Ngữ cảnh theo phiên làm việc & Bộ nhớ dài hạn, lưu vết trạng thái \\
Cơ chế sửa sai & Không tự sửa lỗi & Dựa vào phản hồi của người dùng & Tự phản chiếu (self-reflection), thử lại \\
Vai trò con người & Đặt câu hỏi, đọc kết quả & Thao tác chính, AI gợi ý & Giám sát cấp cao (supervisor / HITL) \\
\bottomrule
\end{tabularx}
\end{table}

\section{Lộ trình bốn giai đoạn tiến hóa công nghệ trong tổ chức}
Hình~\ref{fig:ch2-roadmap} mô tả lộ trình chuyển dịch của AI trong doanh nghiệp. \textbf{Giai đoạn~1} là chatbot thụ
động trả lời câu hỏi thường gặp. \textbf{Giai đoạn~2} là trợ lý Copilot cài sẵn trong công cụ làm việc, tăng năng suất
cá nhân. \textbf{Giai đoạn~3} là tác tử đơn lẻ tự trị, đảm nhận trọn một quy trình. \textbf{Giai đoạn~4} là mạng lưới
nhiều tác tử chuyên biệt phối hợp với nhau trên quy mô toàn doanh nghiệp, McKinsey gọi là \emph{Agentic Mesh}.

\hinhtikz{ch2_roadmap}{Lộ trình tiến hóa công nghệ từ Chatbot thụ động đến mạng lưới tác tử doanh nghiệp (Agentic
Mesh).}{ch2-roadmap}

Mỗi bước tiến kéo theo hai thay đổi: mức tự trị tăng và vai trò của con người dịch chuyển từ \emph{người thao tác} sang
\emph{người giám sát}. Vì vậy các yêu cầu về kiểm soát, nhật ký và phân quyền (Chương~\ref{ch:quantri}) cũng tăng
tương ứng.

\section{Nguyên lý mạng lưới tác tử (Agentic Mesh) theo McKinsey}
McKinsey \cite{mckinsey_agents} mô tả tương lai của tổ chức là một mạng lưới các tác tử chuyên môn hóa sâu, phối hợp qua
các giao thức tương tác giữa tác tử với tác tử (Agent-to-Agent communication) để vận hành những quy trình nghiệp vụ phức
hợp mà không cần con người can thiệp thủ công liên tục. Các đặc điểm chính của mô hình này:
\begin{itemize}
  \item \textbf{Chuyên môn hóa.} Mỗi tác tử phụ trách một miền hẹp (tra cứu đơn hàng, phân tích dữ liệu, kiểm tra tuân
        thủ\dots), dễ kiểm thử và thay thế độc lập.
  \item \textbf{Điều phối.} Một tác tử điều phối (orchestrator) nhận mục tiêu, phân việc và tổng hợp kết quả.
  \item \textbf{Hạ tầng dùng chung.} Các tác tử chia sẻ bộ nhớ, kho công cụ và lớp kiểm toán chung để mọi hành động đều
        truy vết được.
\end{itemize}

\hinhtikz{ch2_mesh}{Mạng lưới tác tử (Agentic Mesh): tác tử điều phối, các tác tử chuyên biệt liên lạc A2A, bộ nhớ --
công cụ dùng chung và lớp kiểm toán.}{ch2-mesh}
